% !TEX root = paper.tex

\section{Formalization} \label{sec:formalization}

The previous sections introduced Verus concepts by example.
This section presents a small formal lambda calculus to make the concepts
from the previous sections more precise.
The goal of this lambda calculus is to serve as a model
to demonstrate particular features and their type safety.
We do not attempt to capture all of the semantics of Rust and Verus,
since formalizing Rust semantics is by itself a large and
challenging problem~\cite{pearce-rust-toplas2021, DBLP:journals/corr/abs-1903-00982, DBLP:journals/pacmpl/0002JKD18}.
Instead, we focus on a small set of topics that are novel to Verus
and are particularly relevant for type safety:

\begin{itemize}
\item \verb`spec`, \verb`proof`, and \verb`exec` functions
\item \verb`spec`, \verb`proof`, and \verb`exec` variables,
  showing how \verb`exec` and \verb`proof` variables are treated linearly,
  while \verb`spec` variables can capture nonlinear snapshots of data
  from \verb`exec` and \verb`proof` variables
\item \verb`spec`, \verb`proof`, and \verb`exec` annotations on datatype fields
\item linear ghost permissions, with read-only borrowing
\item ensuring termination of \verb`spec` and \verb`proof` code,
  particularly in the presence of mutation, recursive types,
  and higher-order features like traits or first-class functions
\ifextended
\item default values in \verb`spec` code,
  particularly for types are are uninhabited in \verb`proof` code and \verb`exec` code
\fi
\end{itemize}

Since this is already a sizable list of topics,
we aggressively minimize other features in our model language.
First, we omit concurrency entirely.
Second, we omit preconditions, postconditions, and verification condition generation,
focusing instead on type safety and termination.
(We believe that verification condition generation could be added
in a style similar to the formalization of Linear Dafny~\cite{DBLP:journals/pacmpl/LiLZCHPH22}.)
Third, our lambda calculus is a mostly-functional language that manipulates values,
rather than an imperative language that mutates values stored in locations.
(This contrasts with more detailed formalizations of Rust centered on
locations~\cite{pearce-rust-toplas2021, DBLP:journals/corr/abs-1903-00982}.)
\ifextended
The model language does, however, include two forms of mutation.
First, the language supports load and store operations that are controlled by linear ghost permissions.
Second, the lambda calculus contains a tiny nonlinear mutable heap
(actually, just a single nonlinear mutable heap location),
for the purpose of demonstrating the type safety of higher-order code in the presence of mutation.
\else
The model language does, however, include mutation,
in the form of load and store operations that are controlled by linear ghost permissions.
(The extended version of this paper~\cite{verus2023extended} also includes a second kind of mutation,
in the form of a tiny nonlinear mutable heap.)
\fi

Since our model language is based on values rather than locations,
it lacks Rust's distinction between a value (e.g. of type \verb`int`)
and a reference to that value (e.g. of type \verb`&int` or \verb`&mut int`).
Nevertheless, we still want to capture some notion of borrowing in order to demonstrate
borrowed linear ghost permissions.
For this, we associate $\linear$ and $\shared$ {\it usages}
with variable typings and expression typings.
The usage ``$\shared$'' represents immutable borrowing (as in \verb`&int`),
which we use for reading permissions.
For simplicity, we omit mutable borrowing,
instead annotating permissions with ``$\linear$''.

We build these usages into the mode system,
in a style similar to Linear Dafny~\cite{DBLP:journals/pacmpl/LiLZCHPH22}
(which in turn built on earlier work by Wadler's ``let!'' feature~\cite{DBLP:conf/ifip2/Wadler90}
and Cogent's purely functional support for borrowing~\cite{cogent}).
We define a mode $\mode$ to be $\spec$, $\prf$, or $\exec$ (see \autoref{fig:formal:syntax}),
with a reflexive, transitive ordering $\exec \modeless \prf \modeless \spec$
and a least upper bound $\mode_1 \modejoin \mode_2$ that
is the least $\mode$ such that
$\mode_1 \modeless \mode$ and
$\mode_2 \modeless \mode$.
We then associate a usage $\usage$ with proof and exec modes,
since proof and exec variables can be linear or borrowed:
\[
\modeu ::= \spec \sep \modu{\prf}{\usage} \sep \modu{\exec}{\usage}
\]
The environment
$\xenv ::= \{x_1 \mapsto \utyp{\modeu_1}{\typ_1}, \ldots, x_n \mapsto \utyp{\modeu_n}{\typ_n}\}$
tracks the mode, usage, and type of each variable.
We refer to a binding $x \mapsto \utyp{\modu{\mode}{\linear}}{\typ}$ as a linear binding,
and we refer to $x \mapsto \utyp{\modu{\mode}{\shared}}{\typ}$
and $x \mapsto \utyp{\spec}{\typ}$ as nonlinear bindings.
We write $\envl\xenv$ to extract just the linear bindings from $\xenv$
and we write $\envn\xenv$ to extract just the nonlinear bindings from $\xenv$
(see \autoref{fig:formal:typingdefs}).

We write $\xenv_1, \xenv_2$ to concatenate two environments together.
For writing typing rules, though,
we often want to split environments in a more sophisticated way then simple concatenation.
In particular, we want to split linear bindings between subexpressions
while sharing nonlinear bindings among subexpressions.
For this, we write $\xenv = \envmerge{\xenv_1}{\xenv_2}$.
For example, in the typing rule for adding two integers (see \autoref{fig:formal:typing}),
the left subexpression gets environment $\xenv_1$ and the right
subexpression gets $\xenv_2$:
\[
\inferrule{
  \ejudgment{\datadecls}{\heaptyp}{\penv_1}{\xenv_1}{\mode}{\laxity}{\expr_1}{\modeu}{\typint}
  \\
  \ejudgment{\datadecls}{\heaptyp}{\penv_2}{\xenv_2}{\mode}{\laxity}{\expr_2}{\modeu}{\typint}
}{
  \ejudgment{\datadecls}{\heaptyp}{\envmerge{\penv_1}{\penv_2}}{\envmerge{\xenv_1}{\xenv_2}}{\mode}{\laxity}{\expr_1 + \expr_2}{\modeu}{\typint}
}
\]
(The other environments can be ignored for now;
\autoref{sec:termination} discusses $\datadecls$ and $\heaptyp$,
and \autoref{sec:permissions} discusses $\penv$ and $\mode$.)

When $\xenv = \envmerge{\xenv_1}{\xenv_2}$,
all nonlinear bindings in $\xenv$ appear in both $\xenv_1$ and $\xenv_2$.
Linear bindings, however are more subtle:
a linear binding in $\xenv$ appears as-is in one of the environments ($\xenv_1$ or $\xenv_2$),
and is demoted to mode $\spec$ in the other environment.
Thus, the environment that didn't get the linear binding can still talk about the variable
in specifications.
For example, if $\xenv$ has a linear binding for $x_2$ and we split $\xenv$
into $\xenv = \envmerge{\xenv_1}{\xenv_2}$,
and $\xenv_1$ receives the linear binding for $x_2$,
then $\xenv_2$ will receive a $\spec$ binding for $x_2$:

$\xenv = \{x_1 \mapsto \utyp{\modu{\exec}{\shared}}{\typ},\: x_2 \mapsto \utyp{\modu{\exec}{\linear}}{\typ}\}$

$\xenv_1 = \{x_1 \mapsto \utyp{\modu{\exec}{\shared}}{\typ},\: x_2 \mapsto \utyp{\modu{\exec}{\linear}}{\typ}\}$

$\xenv_2 = \{x_1 \mapsto \utyp{\modu{\exec}{\shared}}{\typ},\: x_2 \mapsto \utyp{\spec}{\typ}\}$

(See \autoref{fig:formal:typingdefs} for a formal definition of $\envmerge{\xenv_1}{\xenv_2}$.)

Following Linear Dafny's formalization,
our model language allows borrowing by temporarily viewing linear variables as shared
within a lexical scope.
For example, in the sequencing expression $\exprseq{\expr_1}{\expr_2}$
the first expression $\expr_1$ can view a portion of the environment $\xenv_b$ as shared,
and these variables then revert to linear in $\expr_2$:
\[
\inferrule{
  \ejudgment{\datadecls}{\heaptyp}{\penv_1, \envshared{\penv_b}}{\xenv_1, \envshared{\xenv_b}}{\mode}{\laxity}{\expr_1}{\modeu_1}{\typunit}
  \\
  \ejudgment{\datadecls}{\heaptyp}{\penv_2, \envlinear{\penv_b}}{\xenv_2, \envlinear{\xenv_b}}{\mode}{\laxity}{\expr_2}{\modeu_2}{\typ_2}
}{
  \ejudgment{\datadecls}{\heaptyp}{(\envmerge{\penv_1}{\penv_2}), \envlinear{\penv_b}}{(\envmerge{\xenv_1}{\xenv_2}), \envlinear{\xenv_b}}{\mode}{\laxity}{\exprseq{\expr_1}{\expr_2}}{\modeu_2}{\typ_2}
}
\]
Here, the notation $\envlinear{\xenv_b}$ means $\xenv_b$ with all proof/exec bindings made linear,
$\envshared{\xenv_b}$ means $\xenv_b$ with all proof/exec bindings made shared.
(For more detail on this style of borrowing, which was inspired by
Wadler's ``let!'' feature~\cite{DBLP:conf/ifip2/Wadler90}, see~\cite{DBLP:journals/pacmpl/LiLZCHPH22}.)

For simplicity and clarity,
the model language implements a linear type system that prohibits discarding linear resources;
for example, it disallows discarding permissions,
and the only way to deallocate a linear \verb`struct` is to deconstruct it with pattern matching.
(Rust behaves more like an affine type system, allowing dropping of any value.)
Rust includes a \verb`Copy` trait,
implemented by simple types like \verb`bool` and \verb`u64`,
for types that are inherently nonlinear and may be freely copied.
Our model language also includes a judgment $\cjudgment{\datadecls}{\mode}{\typ}$
\ifextended
(see \autoref{fig:formal:additional} for the formal definition)
\else
(see the extended version of this paper~\cite{verus2023extended} for the formal definition)
\fi
to indicate that a type $\typ$ may be copied or dropped,
although, for simplicity, the copies and drops are explicit,
using the expressions $\exprcopy{\expr}$ and $\exprdrop{\expr}$.

Even though the linear type system prohibits copying and dropping linear resources,
it allows arbitrary implicit copying and dropping in specifications.
It defines $\cjudgment{\datadecls}{\spec}{\typ}$ to be true of all types,
so that in $\spec$ mode, code can always use $\exprcopy{\expr}$ and $\exprdrop{\expr}$.
Furthermore, $\spec$ variables can be implicitly copied when splitting
variables among subexpressions using $\envmerge{\xenv_1}{\xenv_2}$.
In particular, since environment splitting creates $\spec$ copies of linear variable bindings,
$\spec$ variables can be used to capture immutable snapshots of mutable linear resources.
For example if variable $x_l$ is bound linearly,
the expression ``$\exprlet{\spec}{x_s}{x_l}{\expr}$''
can make a $\spec$ copy $x_s$ of the linear variable $x_l$ without consuming $x_l$.
Here, $\expr$ can continue to use $x_l$ linearly while simultaneously
keeping the immutable snapshot $x_s$.
This allows specifications to talk about the past state (old snapshots) of linear resources
as well as the current state,
which is useful for specifications that relate old states to new states.

\subsection{Permissions} \label{sec:permissions}

\autoref{sec:ghost:perms} described how linear ghost permissions
allow safe manipulation of low-level pointers.
To model permissions and pointers,
the model language contains a $\typpermi{\typ}$ type representing permission
to read or write a value to pointer $\integer$,
which, for simplicity, is simply an integer constant.
There are three operations on pointers and permissions:

\begin{itemize}
\item $\exprpreadi{\expr_p}$ reads the value stored at pointer $\integer$,
  based on the access granted by permission $\expr_p$
\item $\exprpwritei{\expr_v}{\expr_p}$ writes a new value $\expr_v$ to pointer $\integer$,
  based on the access granted by permission $\expr_p$
\item $\exprpdata{\expr_p}$ takes a $\spec$-mode snapshot of the value currently stored at pointer $\integer$,
  based on the access granted by a $\spec$-mode copy of the permission $\expr_p$
\end{itemize}

\autoref{fig:formal:typingdefs} shows the typing rules for these operations.
Each operation uses the same permission,
but with a different mode.
Writing requires linear access to the permission,
so that no aliased views of the permission can have a stale view of the permission.
Reading, on the other hand, can be performed on borrowed permissions with mode $\shared$.
Finally, $\exprpdata{\expr_p}$ uses the permission with mode $\spec$,
allowing use in specifications.
For simplicity, we omit operations to deallocate permissions or allocate new permissions;
we assume that all permissions are passed in to a program when the program starts and
returned at the end of the program.
However, the typing rule for $\exprpwritei{\expr_v}{\expr_p}$ allows the program to
change the type of a permission, effectively reallocating the memory for a new type.
Thus, the linear handling of permissions is crucial;
if the permissions were not linear, a program could use a stale permission to read
a value memory from memory with an out-of-date type, subverting type safety.
Our type safety theorem (\autoref{sec:soundness}) ensures that this cannot happen.

While $\exprpdata{\expr_p}$ is a ghost-only operation,
$\exprpreadi{\expr_p}$ and $\exprpwritei{\expr_v}{\expr_p}$ perform run-time actions that,
in an implementation, would be compiled to machine code.
Since proofs and specifications are ghost code,
they are not allowed to perform $\exprpreadi{\expr_p}$ and $\exprpwritei{\expr_v}{\expr_p}$ operations.
To enforce this, the typing rules include an access level $\mode$
that limits what operations the code is allowed to perform.
Many operations (such as integer addition) can be performed in any mode,
but $\exprpreadi{\expr_p}$ and $\exprpwritei{\expr_v}{\expr_p}$ can only be performed in $\exec$ mode:
\[
\inferrule{
  \ldots
}{
  \ejudgment{\datadecls}{\heaptyp}{\envmerge{\penv_1}{\penv_2}}{\envmerge{\xenv_1}{\xenv_2}}{\exec}{\laxity}{\exprpwritei{\expr_v}{\expr_p}}{\modu{\prf}{\linear}}{\typ}
}
\]
The bodies of $\exec$ functions are type-checked with access level $\exec$,
and can perform run-time reads and writes,
while the bodies of $\prf$ and $\spec$ functions are type-checked with access level
$\prf$ and $\spec$, and therefore cannot perform run-time reads and writes.

The formal semantics of our mode language use a value $\exprpermi{\val}$
to represent the storage location pointed to by pointer $\integer$,
holding contents $\val$.
Notice that this storage location is just a value,
so it may get passed around from expression to expression,
in and out of functions,
although the type system's linearity ensures that there will never be two inconsistent linear copies
of a permission for the same pointer.
There may, however, be many $\spec$ copies of the permission floating around that
contain snapshots old permission state,
and the code is allowed to execute $\exprpdata{x_s}$ on these snapshots to obtain
the old contents as $\spec$ values.
In fact, it's important that $\exprpdata{x_s}$ return the contents associated with
the snapshot $x_s$,
rather than the most up-to-date value,
because specifications must be deterministic:
they cannot produce different values just because the state has changed.

In order to prove the type safety of our model language,
we have to prove that well-typedness is preserved, including the well-typedness of permission values.
For this, we use an environment $\penv$ that keeps track of whether each storage location $\integer$
is currently linear or borrowed ($\shared$).
We also need to type-check the snapshotted $\spec$ copies of permissions.
For this, we provide a special ``dead-end'' rule that allows stale copies of a permission
to persist as $\spec$-only (see \autoref{fig:formal:typing2});
this is sound because the $\spec$-only permission cannot be coerced
back to a $\shared$ or $\linear$ permission for run-time reads and writes
(hence our description of the $\spec$ copy as a ``dead-end'').

\subsection{Functions and Lifetimes}

Since our model language is a lambda calculus,
it supports functions.
This models both first-order functions,
as shown in the examples in previous section,
and higher-order features.
For example, Verus supports first-class functions in specifications.
Verus also supports simple traits with methods taking a \verb`self` argument;
these simple traits can encode first-class functions.

First-class functions in Rust (called closures in Rust terminology)
are considerably more complicated
than simply-typed lambda calculus functions, though.
First, Rust distinguishes between \verb`Fn`,
which represents functions that may be called many times,
and \verb`FnOnce`,
which represents functions that can only be called once.
(There is also \verb`FnMut`, which we do not model.)
\verb`FnOnce` functions may capture linear variables,
while \verb`Fn` functions cannot.
We define a {\it callability} $\callability ::= \once \sep \many$
to represent this distinction.

Rust first-class functions also have lifetimes associated with them,
so that a function cannot outlive the variables that it captures,
even if these variables are nonlinear (e.g., variables of type \verb`&int`).
Rust lifetimes are quite sophisticated, including parameterization over lifetime variables,
but, for simplicity,
our model language contains just two hard-coded lifetimes
$\lifetime ::= \static \sep \restricted$.
The lifetime $\static$ means that a function may be passed around freely,
because it does not capture any $\shared$ variables,
while the lifetime $\restricted$ means that a function may have captured $\shared$ variables,
and therefore the function cannot be returned past the nearest enclosing borrowing scope.
(Note that this rather strict limitation is only for the model language;
the actual \verus implementation allows Rust's more sophisticated lifetime variables.)
The definition $\textrm{function\_body\_context}$ (see \autoref{fig:formal:typingdefs})
specifies exactly which variables may be captured by the body of a function definition
for each of the four combinations of $\callability$ and $\lifetime$.

Finally, Verus adds yet another dimension to functions: a mode $\mode$
that represents a function being a $\spec$ function, $\prf$ function, or $\exec$ function.
The typing rule for function calls $\exprapp{\expr_f}{\expr_a}$ (\autoref{fig:formal:typing2})
require that the function $\expr_f$'s mode be accessible
according to the current access level,
which means $\mode \modeless \mode_f$ if the current access level is $\mode$
and $\expr_f$ is a function of mode $\mode_f$.

With all of these configuration options,
we can write function definitions $\exprfun{\mode}{\callability}{\lifetime}{x}{\modeu}{\typ}{\expr}$
of type $\typfun{\mode}{\callability}{\lifetime}{\modeu_1}{\typ_1}{\modeu_2}{\typ_2}$.
\autoref{fig:formal:typing2} shows two main rules for assigning function types to function definitions,
one for non-spec functions and one for spec functions.
The latter allows spec functions to capture snapshots of shared variables without worrying about lifetimes;
it does not allow direct capturing of linear variables (since this would effectively discard the linear variable),
but programs can always capture a linear variable indirectly by splitting off
a $\spec$ copy of the linear variable from the surrounding environment and capturing the $\spec$ copy.


The language also allows $\spec$ snapshots of non-spec functions.
Just as snapshots of permissions required a dead-end rule,
functions also require dead-end rules, which bring a slightly annoying technicality.
We could write very simple dead-end rules that just ignore the function's body completely,
and this would be sound,
since a non-spec function snapshotted as a $\spec$ value can never be called,
so the body doesn't matter.
However, our proof of termination in \autoref{sec:soundness} is based on a translation
of our model language into the calculus of inductive constructions (CIC, the logic used by Coq),
and for this translation we need to retain enough of the body to form a well-typed CIC term.
For this, we need to relax the linearity checking in order for the retained body to remain well-typed
in the model language.
\ifextended
Therefore, we parameterize all of the typing rules with a flag $\laxity ::= \lstrict \sep \llax$ that enables ($\lstrict$) or disables
($\llax$) linearity checking, and use $\llax$ for the dead-end function rules.
Note, however, that $\laxity$ is just for assisting the CIC translation,
and does not correspond to anything in the real Verus implementation of type checking and mode checking;
Verus always uses the $\lstrict$ rules.
\else
We write this relaxed checking using the notation $\turnstile_{\llax}$;
the details of this are included in the extended version of this paper~\cite{verus2023extended}.
\fi

\ifextended
\subsubsection{Default values} \label{sec:default}

In Verus, $\spec$ functions are total and correspond closely to SMT total functions,
for the sake of enabling a direct, efficient translation from Verus into SMT queries.
As in Z3, Boogie, and Creusot,
but unlike Dafny, \fstar, and Coq,
Verus $\spec$ functions always return some value of their output type, for all possible inputs.
For example, the division function is uninterpreted when dividing by zero;
it returns some integer, but we can't know which integer.
A Verus function that wraps a division operation is itself well-formed
and returns whatever arbitrary integer the division returns:

\begin{lstlisting}[language=Verus]
#[spec]
fn my_div(i: int, j: int) -> int { i / j }
\end{lstlisting}

On the other hand, non-spec functions are partial:
Verus prohibits division by zero, for example, in $\prf$ and $\exec$ functions.
(If the code above were declared \verb`#[proof]` or \verb`#[exec]`,
it would need a precondition specifying that $j \neq 0$.)
Therefore, it is important to ensure that $\spec$ values don't leak back into non-spec code.
This is especially important for spec functions that return an uninhabited type,
such as Rust's ``Never'' type (written as ``!'' in Rust syntax);
it would be unsound for a value of the Never type to appear as an $\exec$ value.
(For a real-world example of values of uninhabited type causing unsoundness in Dafny,
see~\cite{dafny:unsoundness}.)

To model this, we include the type $\typnever$ in our language,
along with a special default value $\exprbottom$ that has this type in specifications
(and {\it only} in specifications).
To demonstrate that the value doesn't leak into $\prf$ variables or $\exec$ variables,
we include an expression $\exprcrashnever{\expr}$, usable in $\prf$ and $\exec$ modes,
that crashes (fails to step) when given a value of type $\typnever$.
Our type safety proof ensures that this crash never happens.

Building on $\exprbottom$, we can define default values of all well-formed types.
(The details are found in the definition of ``$\defaultsto{\typ}{\val}$'' in \autoref{fig:formal:additional}.)
\fi

\subsection{Termination} \label{sec:termination}

When Verus code is compiled,
all ghost code is erased (not compiled to machine code).
This erasure is sound only if the ghost code always terminates with no side effects.
The access level described in \autoref{sec:permissions} enforces the absence of side effects.
Enforcing termination, though, is more delicate because
mutation and recursive types can often encode nontermination
when combined with higher-order features,
like traits and first-class functions.

To see how this can happen, consider the following two examples, written in OCaml.
The first example creates a mutable reference that holds a function of type \verb`unit -> unit`.
It then stores a new function into the mutable reference
The new function recursively calls itself by reading itself from the reference,
causing an infinite loop:

\begin{lstlisting}
let r: (unit -> unit) ref = ref (fun () -> ()) in
r := (fun () -> !r ());
!r ()
\end{lstlisting}

The second example passes a function to itself as an argument,
using a recursive type \verb`R` to encapsulate the function in a well-typed way.
The function then calls its argument, which means it calls itself,
causing an infinite loop:

\begin{lstlisting}
type r = | R of (r -> unit)
let f (R x) = x (R x) in f (R f)
\end{lstlisting}

\noindent(Note: this can also be encoded directly in Rust as follows:
\begin{lstlisting}[language=Verus]
        trait T { fn f(&self); }
        fn rec<A: T>(x: &A) { x.f(); }
        struct S {}
        impl T for S { fn f(&self) { rec(self); } }
        fn foo() { let s = S {}; s.f(); }
\end{lstlisting}
although this is more complicated.)

Neither of these examples would be caught by \verb`decreases` clauses,
because there are no explicit recursively-defined functions in the code.

\ifextended
To demonstrate that Verus can correctly prohibit these sources of nontermination,
we include two more features in our model language.
First, we add a small heap that consists of a single value $\heapval = \val$
stored in a single heap location, having type $\heaptyp = \typ$.
We also add operations $\exprhread$, $\exprhwrite{\expr}$, and $\exprhdata$
to read, write, and snapshot the heap data.
Using this, we can express nontermination with the following heap value $\heapval$:
\[
\heapval =
  \exprfun{\exec}{\many}{\static}{x}{\modu{\exec}{\linear}}{\typunit}
    {\exprapp{(\exprhread)}{x}}
\]
As in the first OCaml example, this function reads itself from the heap
and then calls itself, causing an infinite loop.
In fact, our model language allows this nontermination for $\exec$ functions.
It prohibits it for $\spec$ and $\prf$ functions,
though, because $\spec$ and $\prf$ functions don't have a sufficent access level $\mode$
to call $\exprhread$, which requires access level $\exec$.
More subtly, they don't have sufficient access to call $\exprhdata$ either,
which could be another source of nontermination ---
unlike the $\exprpdata{\expr_p}$ operation, $\exprhdata$ requires access level $\exec$,
not access level $\spec$.
(This access-level restriction also ensures that specifications are deterministic,
since $\exprhdata$ reads directly from the global heap instead of from a snapshotted permission value.)

Note that Verus currently uses linear permissions rather than a heap,
but other languages like Dafny and \fstar use heaps,
and Linear Dafny uses regions, which have similar properties to heaps.
We include both the heap and linear permissions in the model language
to highlight the distinction between the rules for the two approaches.
In particular, we found the difference between
the $\exprpdata{\expr_p}$ rules and the $\exprhdata$ rules surprising and worth formalizing.

The second feature we add is recursive type definitions
in the form of recursive structs.
\else

We demonstrate that Verus can correctly prohibit these sources of nontermination by including
the following features in the model language:
(i) permissions (\autoref{sec:termination});
(ii) a small heap consisting of one ref cell
(in the extended version of the paper~\cite{verus2023extended});
(iii) recursive types, in the form of recursive structs, described below.

\fi
A recursive struct declaration
$\datadecl ::= \declstruct{\structname}{\structfield{\mode_1}{\typ_1},\ldots,\structfield{\mode_n}{\typ_n}}$
declares a struct $\structname$ with $n$ fields,
each having a mode and a type.
When constructing or destructing structs,
the typing rules join the mode of the fields with the mode of the overall struct value
using the $\modejoin$ operator (see \autoref{fig:formal:typing2}).
This joining ensures, for example, that when reading fields from a $\spec$ snapshot
of an $\exec$ struct value, the result will have mode $\spec$ even if the field mode
is $\prf$ or $\exec$.

\ifextended
To complement recursive structs, the language also includes an option type $\typopt{\typ}$,
allowing interesting recursive types like lists and trees.
The rules for options are straightforward, with one nuance:
the expression for matching on options ($\exprifsome{x}{\expr_1}{\expr_2}{\expr_3}$)
restricts the access level for $\expr_2$ and $\expr_3$ depending on the mode of the option
(see \autoref{fig:formal:typing}).
This prohibits, for example, $\exec$ code from testing whether a $\spec$ option
is Some or None and causing a side effect that depends on the test result.
It does allow $\prf$ functions to test $\spec$ options, though, which is sound and useful in proofs.
\fi

The rules for well-formed struct declarations
\ifextended
\else
(included in the extended version of this paper~\cite{verus2023extended})
\fi
allow recursive structs (although, for simplicity, they disallow mutual recursion).
These rules enforce a standard ``strict positivity'' restriction
(used by Coq, Lean, \fstar, and Dafny).
However, they only require strict positivity in $\spec$ and $\prf$ function types;
non-positive uses are allowed in $\exec$ function types
(unlike in Coq, Lean, and Dafny, where all function types are restricted).

\ifextended
The main judgment for well-formed types has the form
$\trjudgment{\datadecls}{\datadecls_r}{\datadecls_p}{\typ}$ (see \autoref{fig:formal:typing2}).
The environment $\datadecls$ contains all declarations before the current struct
that we're checking,
and $\datadecls_r$ contains the current struct if it is legal
to use recursively in the current position.
The global environment $\datadecls$ is well formed ($\djudgment{\datadecls}$)
if the empty environment is well formed ($\djudgment{\varnothing}$) and
each subsequent struct declaration is well formed:
\[
\inferrule{
  \djudgment{\datadecls}
  \\
  \datadecl = \declstruct{\structname}{\structfield{\mode_1}{\typ_1},\ldots,\structfield{\mode_n}{\typ_n}}
  \\
  \forall 1\le i \le n,\;(\trjudgment{\datadecls}{\varnothing}{\datadecl}{\typ_i}\;\;\textrm{and}\;\;\ljudgment{\mode_i}{\typ_i}{\static})
}{
  \djudgment{\datadecls, \datadecl}
}
\]
For example, the rule for $\spec$ function types makes the struct available
in the return type, but not in the argument type ($\datadecls_r$ is set to $\varnothing$
when checking the argument type):

\[
\inferrule{
  \trjudgment{\datadecls}{\varnothing}{\varnothing}{\typ_1}
  \\
  \trjudgment{\datadecls}{\datadecls_r}{\datadecls_p}{\typ_2}
}{
  \trjudgment{\datadecls}{\datadecls_r}{\datadecls_p}
    {\typfun{\spec}{\many}{\static}{\spec}{\typ_1}{\spec}{\typ_2}}
}
\]

By contrast, the rule for $\exec$ function types allows recursion in the argument type,
so that executable code can still express nontermination through recursive types.
(Note: the two separate environments $\datadecls_r$ and $\datadecls_p$
are used because the rules actually enforce two separate properties:
first, strict positivity and second, that recursive structs have default values.
For the second property, declarations start in $\datadecls_p$ and then shift to $\datadecls_r$.)
\fi

\subsection{Semantics and Type Safety} \label{sec:soundness}

\ifextended
\autoref{fig:formal:evaluation}
\else
The extended version of this paper~\cite{verus2023extended}
\fi
defines evaluation rules
$\heapexp{\heapval}{\expr} \fullstep \heapexp{\heapval'}{\expr'}$
\ifextended
for a heap and expression to take a single step to a new heap and expression.
\else
for an expression to take a single step to a new expression.
\fi
Based on this and the typing rules,
we have proven type preservation,
progress, and ghost-code termination:

\begin{itemize}
\item Preservation:
if $\djudgment\datadecls$
and $\sjudgment{\datadecls}{\heaptyp}{\penv}{\xenv}{\mode}{\heapval}{\expr}{\modeu}{\typ}$
and $\heapexp{\heapval}{\expr} \fullstep \heapexp{\heapval'}{\expr'}$,\newline
then $\sjudgment{\datadecls}{\heaptyp}{\penv}{\xenv}{\mode}{\heapval'}{\expr'}{\modeu}{\typ}$
\item Progress:
if $\djudgment\datadecls$
and $\sjudgment{\datadecls}{\heaptyp}{\penv}{\varnothing}{\mode}{\heapval}{\expr}{\modeu}{\typ}$
and $\expr$ is not a value,\newline
then there is some $\heapexp{\heapval'}{\expr'}$ such that $\heapexp{\heapval}{\expr} \fullstep \heapexp{\heapval'}{\expr'}$.
\item Termination:
if $\mode\in\{\spec, \prf\}$
and $\djudgment\datadecls$
and $\sjudgment{\datadecls}{\heaptyp}{\penv}{\varnothing}{\mode}{\heapval_0}{\expr_0}{\modeu}{\typ}$
then there is no infinite evaluation sequence
$\heapexp{\heapval_0}{\expr_0} \fullstep \heapexp{\heapval_1}{\expr_1} \fullstep \heapexp{\heapval_2}{\expr_2} \fullstep \heapexp{\heapval_3}{\expr_3} \fullstep \ldots$
\end{itemize}

The supplementary material~\cite{verussupplementary} contains proofs of these theorems.
The preservation and progress proofs are straightforward.
The termination proof works by translating the declarations, types, and expressions
into CIC (calculus of inductive constructions) declarations and terms,
and then proving that the CIC declarations and terms are well-typed
and proving a simulation between the CIC reduction steps and the $\heapexp{\heapval}{\expr}$ evaluation steps.
The translation to CIC is fairly simple:
$\spec$ and $\prf$ function types are translated into corresponding
CIC function types, while $\exec$ function types are simply translated into the unit type,
$\exec$ functions are erased completely (translated into the unit value),
\ifextended
$\exprpermi{\val}$ is translated into $\val$, and the heap is erased completely.
\else
and $\exprpermi{\val}$ is translated into $\val$.
\fi








\begin{figure}
\begin{tabular}{lcrl}
variable        & $ x $ & & \\
integer         & $ \integer $ & ::= & $ \ldots, -2, -1, 0, 1, 2, \ldots $ \\
struct name     & $ \structname $ & & \\
usage           & $ \usage $ & ::= & $ \linear \sep \shared $ \\
mode            & $ \mode $ & ::= & $ \spec \sep \prf \sep \exec$ \\
mode + usage\!\!  & $ \modeu $ & ::= & $ \spec \sep \modu{\prf}{\usage} \sep \modu{\exec}{\usage}$ \\
callability     & $ \callability $ & ::= & $ \once \sep \many $ \\
lifetime        & $ \lifetime $ & ::= & $ \static \sep \restricted $ \\
type            & $ \typ $ & ::= & $ \typint \sep \typunit
\ifextended
                             \sep \typnever
\fi
                             \sep \typpermi{\typ} \sep \typopt{\typ} $ \\
                &           & $ \mid $ & $ \structname \sep \typfun{\mode}{\callability}{\lifetime}{\modeu_1}{\typ_1}{\modeu_2}{\typ_2} $ \\
value           & $ \val $ & ::= & $ \integer \sep \exprunit
\ifextended
                             \sep \exprbottom
\fi
                             \sep \exprpermi{v} \sep \exprnone{\typ} \sep \exprsome{\val}{\typ} $ \\
                &           & $ \mid $ & $ \exprstruct{\structname}{\val_1,\ldots,\val_n} \sep \exprfun{\mode}{\callability}{\lifetime}{x}{\modeu}{\typ}{\expr} $ \\
expression      & $ \expr $ & ::= & $ x \sep \integer \sep \expr_1 + \expr_2 \sep \exprunit
\ifextended
                             \sep \exprbottom \sep \exprdefault{\typ} \sep \exprcrashnever{\expr}
\fi
                             $ \\
\ifextended
                &           & $ \mid $ & $ \exprhdata \sep \exprhread \sep \exprhwrite{\expr} $ \\
\fi
                &           & $ \mid $ & $ \exprpermi{v} \sep \exprpdata{\expr_p} \sep \exprpreadi{\expr_p} \sep \exprpwritei{\expr_v}{\expr_p} $ \\
                &           & $ \mid $ & $ \exprdrop{\expr} \sep \exprcopy{\expr} \sep \exprseq{\expr_1}{\expr_2} \sep \exprlet{\mode}{x}{\expr_1}{\expr_2} $ \\
                &           & $ \mid $ & $ \exprnone{\typ} \sep \exprsome{\expr}{\typ} \sep \exprifsome{x}{\expr_1}{\expr_2}{\expr_3} $ \\
                &           & $ \mid $ & $ \exprstruct{\structname}{\expr_1,\ldots,\expr_n} \sep \exprletstruct{\structname}{x_1,\ldots,x_n}{\expr_1}{\expr_2} $ \\
                &           & $ \mid $ & $ \exprfun{\mode}{\callability}{\lifetime}{x}{\modeu}{\typ}{\expr} \sep \exprapp{\expr_1}{\expr_2} $ \\
datatype decl   & $ \datadecl $ & ::= & $ \declstruct{\structname}{\structfield{\mode_1}{\typ_1},\ldots,\structfield{\mode_n}{\typ_n}} $ \\
datatype decls\!\!  & $ \datadecls $ & ::= & $ \datadecl_1,\ldots,\datadecl_n $ \\
\ifextended
heap value      & $ \heapval $ & ::= & $\val$ \\
heap type       & $ \heaptyp $ & ::= & $\typ$ \\
\fi
permission env\!\!\!\!\!\!  & $ \penv $ & ::= & $ \{\integer_1 \mapsto \usage_1, \ldots, \integer_n \mapsto \usage_n\} $ \\
variable env    & $ \xenv $ & ::= & $ \{x_1 \mapsto \utyp{\modeu_1}{\typ_1}, \ldots, x_n \mapsto \utyp{\modeu_n}{\typ_n}\} $ \\
\ifextended
lax checking    & $ \laxity $ & ::= & $ \lstrict \sep \llax $ \\
\fi
\end{tabular}
\caption{Formal Model Language Syntax}
\label{fig:formal:syntax}
\end{figure}













\begin{figure}
\begin{flushleft}
{\small

$\modeof{\spec} = \spec \quad\quad \modeof{\utyp{\mode}{\usage}} = \mode$

$\islinear{\modeu} = \textrm{true}$ iff $\modeu = \modu{\mode}{\linear}$

\begin{tabular}{ll}
$\envn\penv = \{\integer \mapsto \shared \;|\; \integer \mapsto \shared \in \penv\}$
&
$\envn\xenv = \{x \mapsto \utyp{\modeu}{\typ} \;|\; x \mapsto \utyp{\modeu}{\typ} \in \xenv \wedge \neg \islinear{\modeu}\}$
\\
$\envl\penv =  \{\integer \mapsto \linear \;|\; \integer \mapsto \linear \in \penv\}$
&
$\envl\xenv =  \{x \mapsto \utyp{\modeu}{\typ} \;|\; x \mapsto \utyp{\modeu}{\typ} \in \xenv \wedge \islinear{\modeu}\}$
\end{tabular}

\begin{tabular}{ll}
$\envlinear{\penv} = \{\integer \mapsto \linear \;|\; \integer \mapsto \usage \in \penv\}$
&
$\envlinear{\xenv} = \{x \mapsto \utyp{(\modu{\mode}{\linear})}{\typ} \;|\; x \mapsto \utyp{(\modu{\mode}{\usage})}{\typ} \in \xenv\}$
\\
$\envshared{\penv} = \{\integer \mapsto \shared \;|\; \integer \mapsto \usage \in \penv\}$
&
$\envshared{\xenv} = \{x \mapsto \utyp{(\modu{\mode}{\shared})}{\typ} \;|\; x \mapsto \utyp{(\modu{\mode}{\usage})}{\typ} \in \xenv\}$
\\
$ $
&
$\envspec{\xenv} = \{x \mapsto \utyp{\spec}{\typ} \;|\; x \mapsto \utyp{\modeu}{\typ} \in \xenv\}$
\end{tabular}

$\penv = \envmerge{\penv_1}{\penv_2}$ iff $\envl\penv = \envl\penv_1, \envl\penv_2$ and $\envn\penv = \envn\penv_1 = \envn\penv_2$

$\xenv = \envmerge{\xenv_1}{\xenv_2}$ iff $\envl\xenv = \envl\xenv_1, \envl\xenv_2$ and $(\envn\xenv, \envspec{\envl\xenv)} = (\envn\xenv_1, \envspec{\envl\xenv_1)} = (\envn\xenv_2, \envspec{\envl\xenv_2})$

Define $\isunrestricted{\modeu}{\typ}$ to mean: $\modeu \neq (\modu{\mode}{\shared})$
and $\ljudgment{\modeof{\modeu}}{\typ}{\static}$

Define $\isstatic{\xenv}$ to mean: for all $x\mapsto\modu{\modeu}{\typ}\in\xenv$, $\lifetimeof{\typ} = \static$

Define $\nonspecfunctionmodes{m_f}{\modeu_x}{\modeu_b}{\typ_b}$ to mean:
$\isunrestricted{\modeu_b}{\typ_b}$ and
$\mode_f \neq spec$ and
$\mode_f \modeless \modeof{\modeu_x}$ and
$\mode_f \modeless \modeof{\modeu_b}$

Define $\functionbodycontext{\callability}{\lifetime}{\penv}{\xenv}{\penv_b}{\xenv_b}{\usage}$ to mean:
\begin{itemize}
\item if $\callability = \once$ and $\lifetime = \restricted$ then $\penv_b = \penv$ and $\xenv_b = \xenv$
\item if $\callability = \many$ and $\lifetime = \restricted$ then $\penv = \envn\penv$ and $\xenv = \envn\xenv$ and $\penv_b = \penv$ and $\xenv_b = \xenv$
\item if $\callability = \many$ and $\lifetime = \static$ then $\penv = \envn\penv$ and $\penv_b = \varnothing$ and $\xenv = \envn\xenv$ and $\xenv_b = \envspec{\xenv}$
\item if $\callability = \once$ and $\lifetime = \static$ then $\penv_b = \envl\penv$ and $\xenv_b = \envl\xenv, \envspec{\envn\xenv}$ and $\isstatic{\envl\xenv}$
\item if $\callability = \once$ then $\usage = \linear$
\end{itemize}

Define $\lifetimeof{\typ}$ to be:
\begin{itemize}
  \item $\lifetimeof{\typfun{\mode}{\callability}{\lifetime}{\modeu_1}{\typ_1}{\modeu_2}{\typ_2}} = \lifetime$
  \item $\lifetimeof{\typopt{\typ}} = \lifetimeof{\typ}$
  \item $\lifetimeof{\typ} = \static$ for all other $\typ$
\end{itemize}
}
\end{flushleft}
\vspace{-2mm}
\caption{Notation and definitions for type checking}
\label{fig:formal:typingdefs}
\end{figure}











\begin{figure}
\begin{flushleft}{\bf Well-typed expression (main rules)} $\ejudgment{\datadecls}{\heaptyp}{\penv}{\xenv}{\mode}{\laxity}{\expr}{\modeu}{\typ}$\end{flushleft}
{\small

$
\inferrule{
  \mode \modeless \modeof{\modeu_x}
}{
  \ejudgment{\datadecls}{\heaptyp}{\envn\penv}{\envn\xenv, x\mapsto\utyp{\modeu_x}{\typ_x}}{\mode}{\laxity}{x}{\modeu_x}{\typ_x}
}
$
\;\;\;\;\;\;\;\;
$
  \ejudgment{\datadecls}{\heaptyp}{\envn\penv}{\envn\xenv, x\mapsto\utyp{\modu{\mode_x}{\shared}}{\typ_x}}{\mode}{\laxity}{x}{\spec}{\typ_x}
$
\medskip\;\;\;\;\;\;\;\;
$
  \ejudgment{\datadecls}{\heaptyp}{\envn\penv}{\envn\xenv}{\mode}{\laxity}{\integer}{\modeu}{\typint}
$
\medskip\;\;\;\;\;\;\;\;
\inferrule{
  \ejudgment{\datadecls}{\heaptyp}{\penv_1}{\xenv_1}{\mode}{\laxity}{\expr_1}{\modeu}{\typint}
  \\
  \ejudgment{\datadecls}{\heaptyp}{\penv_2}{\xenv_2}{\mode}{\laxity}{\expr_2}{\modeu}{\typint}
}{
  \ejudgment{\datadecls}{\heaptyp}{\envmerge{\penv_1}{\penv_2}}{\envmerge{\xenv_1}{\xenv_2}}{\mode}{\laxity}{\expr_1 + \expr_2}{\modeu}{\typint}
}
\medskip\;\;\;\;\;\;\;\;
$
  \ejudgment{\datadecls}{\heaptyp}{\envn\penv}{\envn\xenv}{\mode}{\laxity}{\exprunit}{\modeu}{\typunit}
$
\medskip\;\;\;\;\;\;\;\;
\ifextended
$
  \ejudgment{\datadecls}{\heaptyp}{\envn\penv}{\envn\xenv}{\mode}{\laxity}{\exprbottom}{\spec}{\typnever}
$
\medskip\;\;\;\;\;\;\;\;
\inferrule{
  \tjudgment{\datadecls}{\typ}
}{
  \ejudgment{\datadecls}{\heaptyp}{\envn\penv}{\envn\xenv}{\mode}{\laxity}{\exprdefault{\typ}}{\spec}{\typ}
}
\medskip\;\;\;\;\;\;\;\;
\inferrule{
  \ejudgment{\datadecls}{\heaptyp}{\penv}{\xenv}{\mode}{\laxity}{\expr}{\modeu}{\typnever}
  \\
  \modeu \neq \spec
}{
  \ejudgment{\datadecls}{\heaptyp}{\penv}{\xenv}{\mode}{\laxity}{\exprcrashnever{\expr}}{\modeu}{\typunit}
}
\medskip\;\;\;\;\;\;\;\;
$
  \ejudgment{\datadecls}{\heaptyp}{\envn\penv}{\envn\xenv}{\exec}{\laxity}{\exprhdata}{\spec}{\heaptyp}
$
\medskip\;\;\;\;\;\;\;\;
$
  \ejudgment{\datadecls}{\heaptyp}{\envn\penv}{\envn\xenv}{\exec}{\laxity}{\exprhread}{\modu{\exec}{\linear}}{\heaptyp}
$
\medskip\;\;\;\;\;\;\;\;
\inferrule{
  \ejudgment{\datadecls}{\heaptyp}{\penv}{\xenv}{\exec}{\laxity}{\expr}{\modu{\exec}{\linear}}{\heaptyp}
}{
  \ejudgment{\datadecls}{\heaptyp}{\penv}{\xenv}{\exec}{\laxity}{\exprhwrite{\expr}}{\modeu}{\typunit}
}
\medskip\;\;\;\;\;\;\;\;
\fi
\inferrule{
  \ejudgment{\datadecls}{\heaptyp}{\envn\penv}{\envn\xenv}{\mode}{\laxity}{\val}{\modu{\exec}{\linear}}{\typ}
  \\
  \cjudgment{\datadecls}{\exec}{\typ}
}{
  \ejudgment{\datadecls}{\heaptyp}{\envn\penv, \integer\mapsto\usage}{\envn\xenv}{\mode}{\laxity}{\exprpermi{\val}}{\modu{\prf}{\usage}}{\typpermi{\typ}}
}
\medskip\;\;\;\;\;\;\;\;
\inferrule{
  \ejudgment{\datadecls}{\heaptyp}{\penv}{\xenv}{\mode}{\laxity}{\expr}{\spec}{\typpermi{\typ}}
}{
  \ejudgment{\datadecls}{\heaptyp}{\penv}{\xenv}{\mode}{\laxity}{\exprpdata{\expr}}{\spec}{\typ}
}
\medskip\;\;\;\;\;\;\;\;
\inferrule{
  \ejudgment{\datadecls}{\heaptyp}{\penv}{\xenv}{\exec}{\laxity}{\expr_p}{\modu{\prf}{\shared}}{\typpermi{\typ}}
}{
  \ejudgment{\datadecls}{\heaptyp}{\penv}{\xenv}{\exec}{\laxity}{\exprpreadi{\expr_p}}{\modu{\exec}{\shared}}{\typ}
}
\medskip\;\;\;\;\;\;\;\;
\inferrule{
  \ejudgment{\datadecls}{\heaptyp}{\penv_1}{\xenv_1}{\exec}{\laxity}{\expr_v}{\modu{\exec}{\linear}}{\typ'}
  \\
  \ejudgment{\datadecls}{\heaptyp}{\penv_2}{\xenv_2}{\exec}{\laxity}{\expr_p}{\modu{\prf}{\linear}}{\typpermi{\typ}}
  \\
  \cjudgment{\datadecls}{\exec}{\typ'}
}{
  \ejudgment{\datadecls}{\heaptyp}{\envmerge{\penv_1}{\penv_2}}{\envmerge{\xenv_1}{\xenv_2}}{\exec}{\laxity}{\exprpwritei{\expr_v}{\expr_p}}{\modu{\prf}{\linear}}{\typ'}
}
\medskip\;\;\;\;\;\;\;\;
\inferrule{
  \ejudgment{\datadecls}{\heaptyp}{\penv}{\xenv}{\mode}{\laxity}{\expr}{\modu{\mode_e}{\linear}}{\typ}
  \\
  \cjudgment{\datadecls}{\mode_e}{\typ}
}{
  \ejudgment{\datadecls}{\heaptyp}{\penv}{\xenv}{\mode}{\laxity}{\exprdrop{\expr}}{\modu{\mode_e}{\shared}}{\typ}
}
\medskip\;\;\;\;\;\;\;\;
\inferrule{
  \ejudgment{\datadecls}{\heaptyp}{\penv}{\xenv}{\mode}{\laxity}{\expr}{\modu{\mode_e}{\shared}}{\typ}
  \\
  \cjudgment{\datadecls}{\mode_e}{\typ}
}{
  \ejudgment{\datadecls}{\heaptyp}{\penv}{\xenv}{\mode}{\laxity}{\exprcopy{\expr}}{\modu{\mode_e}{\linear}}{\typ}
}
\medskip\;\;\;\;\;\;\;\;
\inferrule{
  \ejudgment{\datadecls}{\heaptyp}{\penv_1, \envshared{\penv_b}}{\xenv_1, \envshared{\xenv_b}}{\mode}{\laxity}{\expr_1}{\modeu_1}{\typunit}
  \\
  \ejudgment{\datadecls}{\heaptyp}{\penv_2, \envlinear{\penv_b}}{\xenv_2, \envlinear{\xenv_b}}{\mode}{\laxity}{\expr_2}{\modeu_2}{\typ_2}
}{
  \ejudgment{\datadecls}{\heaptyp}{(\envmerge{\penv_1}{\penv_2}), \envlinear{\penv_b}}{(\envmerge{\xenv_1}{\xenv_2}), \envlinear{\xenv_b}}{\mode}{\laxity}{\exprseq{\expr_1}{\expr_2}}{\modeu_2}{\typ_2}
}
\medskip\;\;\;\;\;\;\;\;
\inferrule{
  \ejudgment{\datadecls}{\heaptyp}{\penv_1, \envshared{\penv_b}}{\xenv_1, \envshared{\xenv_b}}{\mode}{\laxity}{\expr_1}{\modeu_1}{\typ_1}
  \\
  \ejudgment{\datadecls}{\heaptyp}{\penv_2, \envlinear{\penv_b}}{\xenv_2, \envlinear{\xenv_b}, x\mapsto\modu{\modeu_1}{\typ_1}}{\mode}{\laxity}{\expr_2}{\modeu_2}{\typ_2}
  \\
  \isunrestricted{\modeu_1}{\typ_1}\:\textrm{or}\:(\penv_b = \varnothing\:\textrm{and}\:\xenv_b = \varnothing)
  \\
  \ljudgment{\modeof{\modeu_2}}{\typ_2}{\static}
  \\
  \mode_1 = \modeof{\modeu_1}
  \\
  \mode \modeless \mode_1
}{
  \ejudgment{\datadecls}{\heaptyp}{(\envmerge{\penv_1}{\penv_2}), \envlinear{\penv_b}}{(\envmerge{\xenv_1}{\xenv_2}), \envlinear{\xenv_b}}{\mode}{\laxity}
    {\exprlet{\mode_1}{x}{\expr_1}{\expr_2}}{\modeu_2}{\typ_2}
}
\medskip\;\;\;\;\;\;\;\;
\inferrule{
  \tjudgment{\datadecls}{\typ}
}{
  \ejudgment{\datadecls}{\heaptyp}{\envn\penv}{\envn\xenv}{\mode}{\laxity}{\exprnone{\typ}}{\modeu}{\typopt{\typ}}
}
\medskip\;\;\;\;\;\;\;\;
\inferrule{
  \ejudgment{\datadecls}{\heaptyp}{\penv}{\xenv}{\mode}{\laxity}{\expr}{\modeu}{\typ}
}{
  \ejudgment{\datadecls}{\heaptyp}{\penv}{\xenv}{\mode}{\laxity}{\exprsome{\expr}{\typ}}{\modeu}{\typopt{\typ}}
}
\medskip\;\;\;\;\;\;\;\;
\inferrule{
  \ejudgment{\datadecls}{\heaptyp}{\penv_1}{\xenv_1}{\mode}{\laxity}{\expr_1}{\modeu_1}{\typopt{\typ_1}}
  \\
  \ejudgment{\datadecls}{\heaptyp}{\penv_b}{\xenv_b, x\mapsto{\modu{\modeu_1}{\typ_1}}}{\mode_b}{\laxity}{\expr_2}{\modeu_b}{\typ_b}
  \\
  \ejudgment{\datadecls}{\heaptyp}{\penv_b}{\xenv_b}{\mode_b}{\laxity}{\expr_3}{\modeu_b}{\typ_b}
  \\
  \mode \modeless \mode_b
  \\
  (\modeof{\modeu_1} \modeless \mode_b)\:\textrm{or}\:(\modeof{\modeu_1} = \spec\:\textrm{and}\:\mode_b = \prf)
}{
  \ejudgment{\datadecls}{\heaptyp}{\envmerge{\penv_1}{\penv_b}}{\envmerge{\xenv_1}{\xenv_b}}{\mode}{\laxity}{\exprifsome{x}{\expr_1}{\expr_2}{\expr_3}}{\modeu_b}{\typ_b}
}
}
\caption{Type Checking Rules}
\label{fig:formal:typing}
\end{figure}



\begin{figure}
\begin{flushleft}{\bf Well-typed expression (main rules, continued)}\end{flushleft}
{\small
$
\inferrule{
  \datadecls = \ldots, \declstruct{\structname}{\utyp{\mode_1}{\typ_1},\ldots,\utyp{\mode_n}{\typ_n}}, \ldots
  \\
  \forall 1\le i \le n,\;\ejudgment{\datadecls}{\heaptyp}{\penv_i}{\xenv_i}{\mode}{\laxity}{\expr_i}{(\mode_i\modejoin\modeu)}{\typ_i}
}{
  \ejudgment{\datadecls}{\heaptyp}{\envmergen{\penv_1}{\penv_n}}{\envmergen{\xenv_1}{\xenv_n}}{\mode}{\laxity}{\exprstruct{\structname}{\expr_1,\ldots,\expr_n}}{\modeu}{\structname}
}
$
\medskip\;\;\;\;\;\;\;\;
\inferrule{
  \datadecls = \ldots, \declstruct{\structname}{\utyp{\mode_1}{\typ_1},\ldots,\utyp{\mode_n}{\typ_n}}, \ldots
  \\
  \ejudgment{\datadecls}{\heaptyp}{\penv_0}{\xenv_0}{\mode}{\laxity}{\expr_0}{\modeu_0}{\structname}
  \\
  \ejudgment{\datadecls}{\heaptyp}{\penv_b}{\xenv_b, x_1\mapsto\utyp{(\mode_1\modejoin\modeu_0)}{\typ_1}, \ldots, x_n\mapsto\utyp{(\mode_n\modejoin\modeu_0)}{\typ_n}}{\mode}{\laxity}{\expr_b}{\modeu_b}{\typ_b}
  \\
  \ljudgment{\modeof{\modeu_b}}{\typ_b}{\static}
}{
  \ejudgment{\datadecls}{\heaptyp}{\envmerge{\penv_0}{\penv_b}}{\envmerge{\xenv_0}{\xenv_b}}{\mode}{\laxity}{\exprletstruct{\structname}{x_1,\ldots,x_n}{\expr_0}{\expr_b}}{\modeu_b}{\typ_b}
}
\medskip\;\;\;\;\;\;\;\;
\inferrule{
  \ejudgment{\datadecls}{\heaptyp}{\penv_b}{\xenv_b, x\mapsto \utyp{\modeu_x}{\typ_x}}{\mode_f}{\laxity}{\expr_b}{\modeu_b}{\typ_b}
  \\
  \tjudgment{\datadecls}{\typ_x}
  \\
  \functionbodycontext{\callability}{\lifetime}{\penv}{\xenv}{\penv_b}{\xenv_b}{\usage}
  \\
  \nonspecfunctionmodes{\mode_f}{\modeu_x}{\modeu_b}{\typ_b}
}{
  \ejudgment{\datadecls}{\heaptyp}{\penv}{\xenv}{\mode}{\laxity}{(\exprfun{\mode_f}{\callability}{\lifetime}{x}{\modeu_x}{\typ_x}{\expr_b})}{\utyp{\mode_f}{\usage}}{(\typfun{\mode_f}{\callability}{\lifetime}{\modeu_x}{\typ_x}{\modeu_b}{\typ_b})}
}
\medskip\;\;\;\;\;\;\;\;
\inferrule{
  \ejudgment{\datadecls}{\heaptyp}{\envn\penv}{\envn\xenv, x\mapsto \utyp{\spec}{\typ_x}}{\spec}{\laxity}{\expr_b}{\spec}{\typ_b}
  \\
  \tjudgment{\datadecls}{\typ_x}
}{
  \ejudgment{\datadecls}{\heaptyp}{\envn\penv}{\envn\xenv}{\mode}{\laxity}{(\exprfun{\spec}{\many}{\static}{x}{\spec}{\typ_x}{\expr_b})}{\spec}{(\typfun{\spec}{\many}{\static}{\spec}{\typ_x}{\spec}{\typ_b})}
}
\medskip\;\;\;\;\;\;\;\;
\inferrule{
  \callability = \once\Longrightarrow \islinear{\modeu_1}
  \\
  \ejudgment{\datadecls}{\heaptyp}{\penv_1}{\xenv_1}{\mode}{\laxity}{\expr_f}{\modeu_1}{(\typfun{\mode_f}{\callability}{\lifetime}{\modeu_a}{\typ_a}{\modeu_b}{\typ_b})}
  \\
  \ejudgment{\datadecls}{\heaptyp}{\penv_2}{\xenv_2}{\mode}{\laxity}{\expr_a}{\modeu_a}{\typ_a}
  \\
  \modeof{\modeu_1} \modeless \mode_f
  \\
  \mode \modeless \mode_f
}{
  \ejudgment{\datadecls}{\heaptyp}{\envmerge{\penv_1}{\penv_2}}{\envmerge{\xenv_1}{\xenv_2}}{\mode}{\laxity}{\exprapp{\expr_f}{\expr_a}}{\modeu_b}{\typ_b}
}
\ifextended
\medskip\;\;\;\;\;\;\;\;
\inferrule{
  \ejudgment{\datadecls}{\heaptyp}{\penv}{\xenv}{\mode}{\lstrict}{\expr}{\modeu}{\typ}
  \\
  \ejudgment{\datadecls}{\heaptyp}{\varnothing}{\envspec{\xenv}}{\exec}{\lstrict}{\heapval}{\modu{\exec}{\linear}}{\heaptyp}
  \\
  \cjudgment{\datadecls}{\exec}{\heaptyp}
  \\
  \lifetimeof{\heaptyp} = \static
}{
  \sjudgment{\datadecls}{\heaptyp}{\penv}{\xenv}{\mode}{\heapval}{\expr}{\modeu}{\typ}
}
\fi
}

\begin{flushleft}{\bf Well-typed expression (dead-end rules)} $\ejudgment{\datadecls}{\heaptyp}{\penv}{\xenv}{\mode}{\laxity}{\expr}{\modeu}{\typ}$\end{flushleft}
{\small
$
\inferrule{
  \envspec{\xenv} = \envspec{\xenv'}
  \\
  \modeof{\modeu} = \modeof{\modeu'}
  \\
  \ejudgment{\datadecls}{\heaptyp}{\penv'}{\xenv'}{\mode}{\llax}{\expr}{\modeu'}{\typ}
}{
  \ejudgment{\datadecls}{\heaptyp}{\penv}{\xenv}{\mode}{\llax}{\expr}{\modeu}{\typ}
}
$
\medskip\;\;\;\;\;\;\;\;
\inferrule{
  \ejudgment{\datadecls}{\heaptyp}{\envn\penv}{\envn\xenv}{\mode}{\laxity}{\val}{\modeu}{\typ}
}{
  \ejudgment{\datadecls}{\heaptyp}{\envn\penv}{\envn\xenv}{\mode}{\laxity}{\exprpermi{\val}}{\spec}{\typpermi{\typ}}
}
\medskip\;\;\;\;\;\;\;\;
\inferrule{
  \ejudgment{\datadecls}{\heaptyp}{\envn\penv}{\envn\xenv, x\mapsto \utyp{\modeu_x}{\typ_x}}{\mode_f}{\llax}{\expr_b}{\modeu_b}{\typ_b}
  \\
  \tjudgment{\datadecls}{\typ_x}
  \\
  \nonspecfunctionmodes{\mode_f}{\modeu_x}{\modeu_b}{\typ_b}
}{
  \ejudgment{\datadecls}{\heaptyp}{\envn\penv}{\envn\xenv}{\mode}{\laxity}{(\exprfun{\mode_f}{\callability}{\lifetime}{x}{\modeu_x}{\typ_x}{\expr_b})}{\spec}{(\typfun{\mode_f}{\callability}{\lifetime}{\modeu_x}{\typ_x}{\modeu_b}{\typ_b})}
}
\medskip\;\;\;\;\;\;\;\;
\inferrule{
  \ejudgment{\datadecls}{\heaptyp}{\envn\penv}{\envn\xenv, x\mapsto \utyp{\modeu_x}{\typ_x}}{\mode_f}{\llax}{\expr_b}{\modeu_b}{\typ_b}
  \\
  \tjudgment{\datadecls}{\typ_x}
  \\
  \nonspecfunctionmodes{\mode_f}{\modeu_x}{\modeu_b}{\typ_b}
}{
  \ejudgment{\datadecls}{\heaptyp}{\envn\penv}{\envn\xenv}{\mode}{\laxity}{(\exprfun{\mode_f}{\once}{\lifetime}{x}{\modeu_x}{\typ_x}{\expr_b})}{\modu{\mode_f}{\shared}}{(\typfun{\mode_f}{\once}{\lifetime}{\modeu_x}{\typ_x}{\modeu_b}{\typ_b})}
}
}

\ifextended

\begin{flushleft}{\bf Well-formed types} $\trjudgment{\datadecls}{\datadecls_r}{\datadecls_p}{\typ}$ \end{flushleft}
{\small
$
\trjudgment{\datadecls}{\datadecls_r}{\datadecls_p}{\typint}
$
\;\;\;\;\;\;\;\;
$
\trjudgment{\datadecls}{\datadecls_r}{\datadecls_p}{\typunit}
$
\;\;\;\;\;\;\;\;
$
\trjudgment{\datadecls}{\datadecls_r}{\datadecls_p}{\typnever}
$
\;\;\;\;\;\;\;\;
$
\inferrule{
  \trjudgment{\datadecls}{\datadecls_r}{\datadecls_p}{\typ}
}{
  \trjudgment{\datadecls}{\datadecls_r}{\datadecls_p}{\typpermi{\typ}}
}
$
\;\;\;\;\;\;\;\;
\inferrule{
  \trjudgment{\datadecls}{\datadecls_r, \datadecls_p}{\varnothing}{\typ}
}{
  \trjudgment{\datadecls}{\datadecls_r}{\datadecls_p}{\typopt{\typ}}
}
\medskip\;\;\;\;\;\;\;\;
\inferrule{
  \datadecls, \datadecls_r = \ldots, \declstruct{\structname}{\ldots}, \ldots
}{
  \trjudgment{\datadecls}{\datadecls_r}{\datadecls_p}{\structname}
}
\medskip\;\;\;\;\;\;\;\;
\inferrule{
  \trjudgment{\datadecls}{\datadecls_r, \datadecls_p}{\varnothing}{\typ_1}
  \\
  \trjudgment{\datadecls}{\datadecls_r, \datadecls_p}{\varnothing}{\typ_2}
  \\
  \isunrestricted{\modeu_2}{\typ_2}
}{
  \trjudgment{\datadecls}{\datadecls_r}{\datadecls_p}{\typfun{\exec}{\callability}{\lifetime}{\modeu_1}{\typ_1}{\modeu_2}{\typ_2}}
}
\medskip\;\;\;\;\;\;\;\;
\inferrule{
  \prf \modeless \modeof{\modeu_1}
  \\
  \prf \modeless \modeof{\modeu_2}
  \\
  \trjudgment{\datadecls}{\varnothing}{\varnothing}{\typ_1}
  \\
  \trjudgment{\datadecls}{\datadecls_r}{\datadecls_p}{\typ_2}
  \\
  \isunrestricted{\modeu_2}{\typ_2}
}{
  \trjudgment{\datadecls}{\datadecls_r}{\datadecls_p}{\typfun{\prf}{\callability}{\lifetime}{\modeu_1}{\typ_1}{\modeu_2}{\typ_2}}
}
\medskip\;\;\;\;\;\;\;\;
\inferrule{
  \trjudgment{\datadecls}{\varnothing}{\varnothing}{\typ_1}
  \\
  \trjudgment{\datadecls}{\datadecls_r}{\datadecls_p}{\typ_2}
}{
  \trjudgment{\datadecls}{\datadecls_r}{\datadecls_p}{\typfun{\spec}{\many}{\static}{\spec}{\typ_1}{\spec}{\typ_2}}
}

}
\fi

\caption{Type Checking Rules, continued}
\label{fig:formal:typing2}
\end{figure}

\ifextended
\begin{figure}
\begin{flushleft}{\bf Evaluation context} $\efill{\expr}$\end{flushleft}
{\small

\begin{tabular}{lrrl}
& $ \ectx $ & ::= & $
  \ehole
  \sep \ectx_1 + \expr_2
  \sep \val_1 + \ectx_2
  \sep \exprcrashnever{\ectx}
  \sep \exprhwrite{\ectx}
  \sep \exprpdata{\ectx}
  $ \\
  & & $ \mid $ & $
  \exprpreadi{\ectx}
  \sep \exprpwritei{\ectx}{\expr_p}
  \sep \exprpwritei{\val}{\ectx}
  \sep \exprdrop{\ectx}
  \sep \exprcopy{\ectx}
  $ \\
  & & $ \mid $ & $
  \exprseq{\ectx_1}{\expr_2}
  \sep \exprlet{\mode}{x}{\ectx_1}{\expr_2}
  \sep \exprsome{\ectx}{\typ}
  \sep \exprifsome{x}{\ectx}{\expr_2}{\expr_3}
  $ \\
  & & $ \mid $ & $
  \exprstruct{\structname}{\val_1,\ldots,\val_i,\ectx_j,\expr_k,\ldots,\expr_n}
  \sep \exprletstruct{\structname}{x_1,\ldots,x_n}{\ectx_0}{\expr_b}
  \sep \exprapp{\ectx_1}{\expr_2}
  \sep \exprapp{\val_1}{\ectx_2}
  $ \\
\end{tabular}

}
\begin{flushleft}{\bf Evaluation rules} $(\heapval, \expr) \fullstep (\heapval', \expr')$ (implicitly in a context $\datadecls$)\end{flushleft}
{\small

\[
\inferrule{
  \expr\step\expr'
}{
  (\heapval, \expr) \step (\heapval, \expr')
}
\;\;\;\;\;\;\;\;
\inferrule{
  (\heapval, \expr) \step (\heapval', \expr')
}{
  (\heapval, \efill{\expr})\fullstep(\heapval', \efill{\expr'})
}
\]

\[
\inferrule{
  i_3 = \textrm{sum of } i_1, i_2
}{
  i_1 + i_2 \step i_3
}
\;\;\;\;\;\;\;\;
\inferrule{
  \defjudgment{\datadecls}{\typ}{\val}
}{
  \exprdefault{\typ} \step \val
}
\]

\[
(\heapval, \exprhdata) \step (\heapval, \heapval)
\;\;\;\;\;\;\;\;
(\heapval, \exprhread) \step (\heapval, \heapval)
\;\;\;\;\;\;\;\;
(\heapval, \exprhwrite{\val}) \step (\val, \exprunit)
\]

\[
\exprpdata{\exprpermi{\val}} \step \val
\;\;\;\;\;\;\;\;
\exprpreadi{\exprpermi{\val}} \step \val
\]

\[
\exprpwritei{\val'}{\exprpermi{\val}} \step \exprpermi{\val'}
\]

\[
\exprdrop{\val} \step ()
\;\;\;\;\;\;\;\;
\exprcopy{\val} \step \val
\;\;\;\;\;\;\;\;
\exprseq{\exprunit}{\expr_2} \step \expr_2
\;\;\;\;\;\;\;\;
\exprlet{\mode}{x}{\val_1}{\expr_2} \step \subst{\expr_2}{x}{\val_1}
\]

\[
\exprifsome{x}{\exprnone{\typ}}{\expr_2}{\expr_3} \step \expr_3
\]

\[
\exprifsome{x}{\exprsome{\val}{\typ}}{\expr_2}{\expr_3} \step \subst{\expr_2}{x}{\val}
\]

\[
\exprletstruct{\structname}{x_1,\ldots,x_n}{\exprstruct{\structname}{\val_1,\ldots,\val_n}}{\expr_b}
\step \msubst{\expr_b}{x_1}{\val_1}{x_n}{\val_n}
\]

\[
\exprapp{(\exprfun{\mode}{\callability}{\lifetime}{x}{\modeu_x}{\typ_x}{\expr_b})}{\val_x} \step \subst{\expr_b}{x}{\val_x}
\]

$ $

(Note: $\exprcrashnever{\exprbottom}$ does not step.  By not stepping, it ``crashes''.)

}

\caption{Evaluation Rules}
\label{fig:formal:evaluation}
\end{figure}


\begin{figure}
\begin{flushleft}{\bf Default values} $\defjudgment{\datadecls}{\typ}{\val}$\end{flushleft}
{\small

\[
\defjudgment{\datadecls}{\typint}{0}
\;\;\;\;\;\;\;\;
\defjudgment{\datadecls}{\typunit}{\exprunit}
\;\;\;\;\;\;\;\;
\defjudgment{\datadecls}{\typnever}{\exprbottom}
\]

\[
\inferrule{
  \defjudgment{\datadecls}{\typ}{\val}
}{
  \defjudgment{\datadecls}{\typpermi{\typ}}{\exprpermi{\val}}
}
\;\;\;\;\;\;\;\;
\defjudgment{\datadecls}{\typopt{\typ}}{\exprnone{\typ}}
\]

\[
\inferrule{
  \datadecls = \ldots, \declstruct{\structname}{\utyp{\mode_1}{\typ_1},\ldots,\utyp{\mode_n}{\typ_n}}, \ldots
  \\
  \defjudgment{\datadecls}{\typ_1}{\val_1}
  \\
  \ldots
  \\
  \defjudgment{\datadecls}{\typ_n}{\val_n}
}{
  \defjudgment{\datadecls}{\structname}{\exprstruct{\structname}{\val_1,\ldots,\val_n}}
}
\]

\[
\defjudgment{\datadecls}
  {(\typfun{\mode}{\callability}{\lifetime}{\modeu_1}{\typ_1}{\modeu_2}{\typ_2})}
  {(\exprfun{\mode}{\callability}{\lifetime}{x}{\modeu_1}{\typ_1}{\exprdefault{\typ_2}})}
\]

}
\begin{flushleft}{\bf Copyable types} $\cjudgment{\datadecls}{\mode}{\typ}$\end{flushleft}
{\small

\[
\cjudgment{\datadecls}{\spec}{\typ}
\;\;\;\;\;\;\;\;
\cjudgment{\datadecls}{\mode}{\typint}
\;\;\;\;\;\;\;\;
\cjudgment{\datadecls}{\mode}{\typunit}
\;\;\;\;\;\;\;\;
\cjudgment{\datadecls}{\mode}{\typnever}
\]

\[
\inferrule{
  \datadecls = \ldots, \declstruct{\structname}{\utyp{\mode_1}{\typ_1},\ldots,\utyp{\mode_n}{\typ_n}}, \ldots
  \\
  \cjudgment{\datadecls}{\mode_1}{\typ_1}
  \\
  \ldots
  \\
  \cjudgment{\datadecls}{\mode_n}{\typ_n}
}{
  \cjudgment{\datadecls}{\mode}{\structname}
}
\]

\[
\inferrule{
  \cjudgment{\datadecls}{\mode}{\typ}
}{
  \cjudgment{\datadecls}{\mode}{\typopt{\typ}}
}
\;\;\;\;\;\;\;\;
  \cjudgment{\datadecls}{\mode}{(\typfun{\mode_f}{\many}{\lifetime}{\modeu_1}{\typ_1}{\modeu_2}{\typ_2})}
\]

}
\caption{Additional Rules}
\label{fig:formal:additional}
\end{figure}

\fi

